% SOURCE_AUTHORITY: sga5_fr_workpass.tex, lines 14145--14192 (LF-normalized unit).
% SOURCE_SHA256: 791F4EFFC5E02832D5D77ED03518C8156D6F07E4C8238B03545DB93D883FBB28
\begin{statement}{Proposição 7.2}
Sejam $X$ um esquema suave de dimensão $1$, $F$ construtível e
\[
(u,f):(X,F)\longrightarrow (X,F)
\]
um endomorfismo tal que o morfismo de feixes
\[
u:f^*(F)\longrightarrow F
\]
seja um isomorfismo. Nestas condições existem $U,U',G,f',\varphi$ que satisfazem as condições de 5.1, 6.1, 6.2.
\end{statement}

\begin{prooftext}[Demonstração]
Seja $U$ o maior aberto de $X$ tal que $F$ seja localmente constante em $U$, e denotemos por $Y$ seu complemento. O conjunto $Y$ é estável por $f$:
\[
f(Y)\subset Y.
\]
Com efeito, suponhamos, por redução ao absurdo, que exista $y\in Y$ tal que $f(y)\in U$. O feixe $f^*(F)$ é então localmente constante em uma vizinhança de $y$ e, usando o isomorfismo $u$, resulta que $F$ é localmente constante em uma vizinhança de $y$, o que é absurdo.

Se $f$ é constante, a demonstração da proposição é trivial. Suporemos, pois, que $f$ não é constante. Se $\eta$ é o ponto genérico de $X$, ter-se-á então $f(\eta)=\eta$.

Sejam $i_\eta:\Spec(k(\eta))\hookrightarrow X$ a inclusão canônica, $F_\eta=i_\eta^*(F)$ e
\[
u_\eta:f_\eta^*(F_\eta)\longrightarrow F_\eta
\]
o morfismo, com efeito um isomorfismo, de feixes induzido por $u$, sendo $f_\eta$, naturalmente, a restrição de $f$ a $\Spec(k(\eta))$.

Traduzindo tudo em termos de $H$-módulos:
\[
H=\Gal(k(\eta)_s/k(\eta)),\qquad k(\eta)_s=\text{um fecho separável de }k(\eta),
\]
obtêm-se os isomorfismos
\[
\begin{tikzcd}[column sep=large]
H \arrow[r,"\Psi"] & H
\end{tikzcd}
\qquad
\begin{tikzcd}[column sep=large]
F_{\bar\eta} \arrow[r,"u_{\bar\eta}"] & F_{\bar\eta},
\end{tikzcd}
\]
onde $F_{\bar\eta}$ tem uma estrutura de $H$-módulo e verifica-se a seguinte relação:
\[
u_{\bar\eta}(gx)=\Psi(g)\,u_{\bar\eta} x,
\qquad g\in H,
\quad x\in F_{\bar\eta}.
\]
Se $H_1$ é o maior subgrupo distinguido de $H$ de índice finito tal que $F_{\bar\eta}^{H_1}=F_{\bar\eta}$, temos $\Psi(H_1)=H_1$. Existe tal subgrupo porque o feixe $F_\eta$ é construtível.
